\documentclass[10pt,a4paper]{article}
\usepackage{fontspec,amsmath,amssymb,graphicx}
\usepackage[margin=18mm]{geometry}
\setmainfont{Times New Roman}
\setlength{\parindent}{0pt}
\setlength{\parskip}{5pt}
\allowdisplaybreaks
\sloppy
\emergencystretch=2em
\begin{document}
\section*{BT\_5\_1\_3\_03 — Lời giải bổ sung}
Với hiệu $w$ của hai nghiệm, $\int_\Omega|\nabla w|^2=\int_{\partial\Omega}w\partial_\nu w=0$, nên $w$ hằng; vì vậy nếu có một nghiệm thì mọi nghiệm là $u+C$.

Điều kiện tương thích (dùng đạo hàm tọa độ trong đề) là
\[\int_0^L(\varphi_1-\varphi_0)dx+\int_0^M(\psi_1-\psi_0)dy=0.\]
Đủ tồn tại nghiệm yếu Neumann, duy nhất sau khi ấn định trung bình; để có nghiệm cổ điển đến góc còn cần độ trơn và tương thích góc.

(iii) Dữ liệu cho $\int(\varphi_1-\varphi_0)=-1/2$, hai cạnh kia 0. Do đó \textbf{không có nghiệm}; không được triển khai Fourier bỏ qua mode 0.\end{document}